\chapter{Irritable Bowel Syndrome (IBS)}

\section*{Definition and Overview}
Irritable bowel syndrome (IBS) is a common disorder of the gut-brain interaction that causes abdominal pain, bloating, and changes in bowel habit, including diarrhoea, constipation, or both. It is a functional disorder, meaning the bowel looks normal on tests but does not work normally, and it is not caused by inflammation, infection, or cancer. IBS is a chronic condition with symptoms that flare and settle, and although it can significantly affect quality of life, it does not damage the bowel or shorten life. Management is individualised and includes dietary change, stress management, and, where needed, medicines and psychological therapies.

\section*{Epidemiology}
IBS affects around one in ten Australians and is one of the most common reasons for gastroenterology referral. It is about twice as common in women as in men and most often begins in late adolescence or early adulthood. Many people with IBS never consult a doctor, and of those who do, most are managed in primary care. IBS commonly coexists with other conditions, including anxiety, depression, fibromyalgia, chronic fatigue, and chronic pelvic pain, and it accounts for a large amount of time off work and reduced productivity. Symptoms often begin or worsen after a gut infection or a period of stress.

\section*{Aetiology and Pathophysiology}
IBS is understood through the biopsychosocial model: the gut and brain communicate continuously, and in IBS this communication is disturbed. The bowel may be hypersensitive, so that normal gas and movement are perceived as pain, and its motility may be abnormal, producing diarrhoea or constipation. The gut microbiome, the community of bacteria in the bowel, differs in people with IBS, and gut infections, food intolerances, and stress can trigger or perpetuate symptoms. Anxiety and depression are common and worsen symptoms through the brain-gut axis, while symptoms themselves generate further distress. There is no single cause, and the relative contribution of each factor varies from person to person, which is why treatment is individualised.

\section*{Clinical Features}
\begin{itemize}
    \item Abdominal pain or cramping, often related to bowel movements and eased by passing stool
    \item Bloating and a feeling of distension
    \item Diarrhoea, constipation, or alternating between the two
    \item Mucus in the stool
    \item A feeling of incomplete emptying after opening the bowels
    \item Urgency to open the bowels
    \item Symptoms that worsen with stress, certain foods, or during periods
    \item No bleeding, weight loss, or fever, which would point to other conditions
\end{itemize}

\section*{Differential Diagnosis}
Because IBS has no diagnostic test, other causes of similar symptoms must be considered.
\begin{itemize}
    \item \textbf{Coeliac disease:} gluten intolerance with similar symptoms
    \item \textbf{Inflammatory bowel disease:} Crohn's disease and ulcerative colitis
    \item \textbf{Infections:} gastroenteritis and parasites
    \item \textbf{Food intolerances:} lactose and other carbohydrate intolerances
    \item \textbf{Bowel cancer:} especially in older people with new symptoms or bleeding
    \item \textbf{Diverticular disease and microscopic colitis:} other bowel conditions
    \item \textbf{Endometriosis:} in women, pelvic symptoms can overlap
    \item \textbf{Thyroid disease:} affecting bowel habit
\end{itemize}

\section*{When to Seek Medical Care}
Anyone with persistent bowel symptoms should be assessed, especially if there is blood in the stool, weight loss, fever, or symptoms beginning after age 50. People with known IBS should seek review if symptoms change, worsen, or do not respond to management, or if new warning features develop. Pregnancy, new medicines, and major stress can alter symptoms and warrant review.

\textbf{Call 000 (or local emergency services) for:}
\begin{itemize}
    \item Severe abdominal pain with a rigid, distended abdomen
    \item Heavy rectal bleeding or blood with faintness and dizziness
    \item Repeated vomiting or inability to keep fluids down
    \item High fever with abdominal pain
    \item Sudden severe symptoms suggesting a bowel obstruction
\end{itemize}

\section*{Diagnosis and Investigations}
\begin{itemize}
    \item \textbf{Clinical criteria:} symptoms consistent with IBS after a careful history and examination
    \item \textbf{Blood tests:} full blood count, coeliac serology, thyroid function, and inflammatory markers
    \item \textbf{Stool tests:} faecal calprotectin and cultures to exclude inflammation and infection
    \item \textbf{Bowel habit diary:} tracking symptoms, foods, and stress
    \item \textbf{Colonoscopy:} recommended in older people or those with warning features
    \item \textbf{Lactose and other food intolerance testing:} where relevant
    \item \textbf{Screening for anxiety and depression:} common and treatable contributors
\end{itemize}

\section*{Management}
\begin{itemize}
    \item \textbf{Lifestyle and diet:} regular meals, adequate fluids, and fibre adjustment according to symptoms
    \item \textbf{Low FODMAP diet:} a structured diet that reduces fermentable carbohydrates, guided by a dietitian
    \item \textbf{Avoiding triggers:} identifying personal food and stress triggers
    \item \textbf{Stress management:} relaxation, mindfulness, and pacing
    \item \textbf{Psychological therapies:} cognitive behavioural therapy and gut-directed hypnotherapy
    \item \textbf{Medicines:} antispasmodics for pain, laxatives or antidiarrhoeals for bowel habit
    \item \textbf{Newer medicines:} low-dose antidepressants and specific IBS medicines for severe symptoms
    \item \textbf{Probiotics:} some evidence of benefit, with strain-specific advice
    \item \textbf{Regular review:} adjusting treatment to the symptom pattern
\end{itemize}

\section*{Complications}
\begin{itemize}
    \item Reduced quality of life and work productivity
    \item Sleep disturbance and fatigue
    \item Anxiety and depression
    \item Avoidance of social activities and travel
    \item Dietary restriction leading to nutritional inadequacy
    \item Unnecessary investigations and treatments when symptoms are misunderstood
    \item Interference with daily activities and relationships
\end{itemize}

\section*{Prognosis and Outcomes}
IBS is a chronic condition, but most people improve significantly with the right combination of dietary, lifestyle, and psychological strategies, and many achieve good control. Symptoms fluctuate, often with stress and diet, and periods of remission are common. IBS does not progress to cancer or inflammatory bowel disease. With individualised management and support, most people with IBS live full and active lives, and severe disability is uncommon.

\section*{Prevention and Self-Care}
\begin{itemize}
    \item Eat regular meals and avoid skipping meals
    \item Drink enough fluid and limit alcohol and caffeine
    \item Adjust fibre according to your symptoms, with dietetic advice
    \item Consider a low FODMAP diet under the guidance of a dietitian
    \item Keep a symptom diary to identify triggers
    \item Manage stress with relaxation, exercise, and sleep
    \item Seek help for anxiety and low mood
    \item Avoid over-the-counter remedies without advice
    \item Attend regular reviews and report new or changing symptoms
\end{itemize}

\section*{Sources and Further Reading}
The following reputable sources provide further information for healthcare students and patients seeking reliable, current advice about irritable bowel syndrome.
\begin{itemize}
    \item Gastroenterological Society of Australia. \textit{IBS patient information.} https://www.gesa.org.au/
    \item Healthdirect Australia. \textit{Irritable bowel syndrome.} https://www.healthdirect.gov.au/
    \item Dietitians Australia. \textit{Low FODMAP diet and IBS.} https://dietitiansaustralia.org.au/
    \item Coeliac Australia. \textit{IBS and coeliac disease information.} https://www.coeliac.org.au/
    \item NHS. \textit{Irritable bowel syndrome.} https://www.nhs.uk/conditions/irritable-bowel-syndrome/
    \item International Foundation for GI Disorders. \textit{IBS education.} https://aboutibs.org/
\end{itemize}
